% part: intuitionistic-logic
% chapter: propositions-as-types
% section: introduction

\documentclass[../../../include/open-logic-section]{subfiles}

\begin{document}

\olfileid{int}{pty}{int}

\olsection{مقدمه}

حساب لامبدا مدلی ساده از محاسبه است که بر ترم‌های ساخته‌شده از
متغیرها عمل می‌کند. اگر $N$ و $M$ ترم باشند، $(NM)$ نیز ترم است
(«اعمال~$N$ بر~$M$»). اگر $N$ ترم باشد، $\lambd[x][N]$ نیز ترم است.
اگر $N$ شامل متغیر~$x$ باشد، وقوع‌های متناظر~$x$ در $\lambd[x][N]$
به‌وسیلهٔ عملگر $\lambd[x]$ مقید می‌شوند و دیگر آزاد نیستند. می‌گوییم
ترمی به شکل $(\lambd[x][N]M)$ با یک گام $\beta$ به
$\Subst{N}{M}{x}$ منقبض می‌شود؛ یعنی به حاصل جایگزینی هر وقوع آزاد
$x$ در~$N$ با~$M$. اگر $N'$ با انقباض $\beta$ یک زیرترم از~$N$ حاصل
شود، می‌گوییم $N$ در یک گام $\beta$ به~$N'$ کاهش می‌یابد و می‌نویسیم
$N \redone N'$. محاسبه‌ای در حساب~$\lambd$ دنباله‌ای چون
$N \redone N_1 \redone N_2 \redone \dots$ است. اگر این دنباله به
ترم~$N_k$ ختم شود و هیچ زیرترم~$N_k$ قابل انقباض~$\beta$ نباشد،
$N_k$ را \emph{بهنجار} یا یک صورت بهنجار~$N$ می‌نامیم. محاسبه در
حساب~$\lambd$ را چنین دنباله‌هایی در نظر بگیرید. محاسبه هنگامی متوقف
می‌شود که به صورت بهنجار برسد، و صورت بهنجار نتیجهٔ محاسبه است. این
مدل سادهٔ محاسبه تورینگ‌کامل است: هر تابع بازگشتی جزئی را می‌توان با
ترمی در حساب~$\lambd$ محاسبه کرد.

هاسکل کِری و ویلیام هاوارد مستقل از یکدیگر شباهتی نزدیک میان استنتاج
طبیعی و حساب~$\lambd$ کشف کردند. به‌طور شهودی، ترم
$\lambd[x][N]$ یک تابع است: $x$ آرگومان آن است و $N$ می‌گوید مقدار را
با داشتن~$x$ چگونه محاسبه کنیم. ترم $(\lambd[x][N]M)$ اعمال تابع
$\lambd[x][N]$ بر آرگومان~$M$ است و انقباض~$\beta$ می‌گوید آن را
چگونه ارزیابی کنیم: $x$ را در~$N$ با~$M$ جایگزین کنید و ارزیابی ترم
حاصل را ادامه دهید. اکنون این‌ها را تابع‌هایی با دامنه و برد ثابت
در نظر بگیرید. اگر دامنهٔ $\lambd[x][N]$ برابر~$A$ و بردش~$B$ باشد،
باید متغیر~$x$ را تنها پذیرای آرگومان‌هایی از~$A$ و $N$ را تولیدکنندهٔ
مقادیر~$B$ بدانیم. در نتیجه، $(\lambd[x][N]M)$ تنها هنگامی معنادار
است که $\lambd[x][N]$ تابعی $A \to B$ و $M \in A$ باشد. افزودن این
اطلاعات به حساب~$\lambd$، \emph{حساب $\lambd$ نوع‌دار} را به دست
می‌دهد. در حساب~$\lambd$ نوع‌دار هر عبارت $(\lambd[x][N]M)$ مجاز
نیست؛ تنها عبارت‌هایی مجازند که «نوع» $\lambd[x][N]$ و $M$ با هم
سازگار باشد: $\lambd[x][N]$ نوع $A \to B$ و $M$ نوع~$A$ داشته باشد.
نوع $\lambd[x][N]$ تنها هنگامی $A \to B$ است که نوع~$x$ برابر~$A$ و
نوع~$N$ برابر~$B$ باشد. در حالت کلی، هر عبارت $(M_1M_2)$ مجاز نیست.
تنها ترم‌هایی از این شکل مجازند که $M_1$ نوع $A \to B$ و $M_2$ نوع
~$A$ داشته باشد؛ در این صورت $(M_1M_2)$ نوع~$B$ دارد.

اکنون این قواعد نوع‌دهی آشکارا با !!{derivation}s استنتاج طبیعی
متناظرند؛ نوع‌ها با !!{formula}s و خود !!{derivation}s با
ترم‌های~$\lambd$ بازنمایی می‌شوند. برای نمونه، !!a{derivation} برای
$!A \lif !B$ با ترم $\lambd[\typeof{x}{!A}][N]$ متناظر است که نوع
تابعی $!A \lif !B$ دارد. قواعد استنتاج طبیعی با قواعد نوع‌دهی در
حساب لامبدای نوع‌دار متناظرند؛ برای نمونه:
\begin{prooftree}
  \AxiomC{$\Discharge{!A}{x}$}
  \DeduceC{$!B$}
  \DischargeRule{\Intro\lif}{x}
  \UnaryInfC{$!A \lif !B$}
  \DisplayProof\bottomAlignProof
  \quad\text{متناظر است با}\quad
  \Axiom$x: !A \fCenter N: !B$
  \UnaryInf$ \fCenter \lambd[\typeof{x}{!A}][N] : !A \lif !B$
\end{prooftree}
قاعدهٔ سمت راست می‌گوید اگر $x$ نوع~$!A$ و $N$ نوع~$!B$ داشته باشد،
آنگاه $\lambd[\typeof{x}{!A}][N]$ نوع~$!A \lif !B$ دارد.

قاعدهٔ $\Elim\lif$ با قاعدهٔ نوع‌دهی برای ترم‌های کاربرد متناظر است:
\[
  \AxiomC{$!A \lif !B$}
  \AxiomC{$!A$}
  \RightLabel{\Elim\lif}
  \BinaryInfC{$!B$}
  \DisplayProof
  \quad\text{متناظر است با}\quad
  \Axiom$\fCenter M_1: !A \lif !B$
  \Axiom$\fCenter M_2: !A$
  \BinaryInf$ \fCenter (M_1M_2) : !B$
  \DisplayProof
\]
اگر یک قاعدهٔ \Intro\lif{} بی‌درنگ در پی یک قاعدهٔ \Elim\lif{}
بیاید، می‌توان !!{derivation} را ساده کرد:
\begin{gather*}
  \AxiomC{$\Discharge{!A}{x}$}
  \DeduceC{$!B$}
  \DischargeRule{\Intro\lif}{x}
  \UnaryInfC{$!A \lif !B$}
  \AxiomC{}
  \DeduceC{$!A$}
  \RightLabel{\Elim\lif}
  \BinaryInfC{$!B$}
  \DisplayProof
  \quad\redone\quad
  \AxiomC{}
  \DeduceC{$!A$}
  \DeduceC{$!B$}
  \DisplayProof
\end{gather*}
که با انقباض~$\beta$ ترم‌های لامبدا متناظر است:
\[
(\lambd[\typeof{x}{!A}][N]M) \quad\redone\quad
\Subst{N}{M}{x}.
\]

تناظرهای مشابهی میان قواعد $\land$ و نوع‌های «ضرب»، و میان قواعد
$\lor$ و نوع‌های «جمع» برقرارند.

این تناظر میان ترم‌های حساب لامبدای ساده‌نوع‌دار و
!!{derivation}s استنتاج طبیعی، تناظر «کِری--هاوارد»، «برهان‌ها
به‌مثابهٔ برنامه‌ها» یا «گزاره‌ها به‌مثابهٔ نوع‌ها» نامیده می‌شود.
علاوه بر تناظر !!{formula}s (گزاره‌ها) با نوع‌ها و برهان‌ها با ترم‌ها،
می‌توان تناظرها را چنین خلاصه کرد:
\begin{center}
  \begin{tabular}{c c}
    منطق & برنامه \\
    \hline
    گزاره/!!{formula} & نوع \\
    برهان & ترم \\
    فرض & متغیر \\
    فرض مرخص‌شده & متغیر مقید \\
    فرض مرخص‌نشده & متغیر آزاد \\
    $!A \lif !B$ & نوع تابعی \\
    $!A \land !B$ & نوع ضرب \\
    $!A \lor !B$ & نوع جمع \\
    $\lfalse$ & نوع تهی \\
    \hline
  \end{tabular}
\end{center}

تناظر کِری--هاوارد یکی از سنگ‌بناهای دستیارهای خودکار اثبات و
نوع‌سنج‌های برنامه است، زیرا بررسی برهانی که گواه یک گزاره است---همان
کاری که در بالا کردیم---هم‌ارز بررسی این است که آیا برنامه (ترم) نوع
اعلام‌شده را دارد یا نه.
\end{document}
